\chapter{Pruebas Funcionales}
Las pruebas funcionales son escenciales en el desarrollo de software, ya que permiten
mostrar y evidenciar el correcto funcionamiento de la aplicaci\'on. En este cap\'itulo 
se realizar\'an distintas pruebas con el objetivo de verificar si todos los requisitos, 
casos de uso y objetivos de esta tesis, se lograron resolver. 

\section{Pruebas de Flujo Administrador}
En este apartado se llevar\'an a cabo pruebas de flujo para comprobar el correcto funcionamiento
de la aplicaci\'on, mostrando im\'agenes reales. 

\subsection{Solicitud de Registro}
Con esta prueba se busca mostrar el correcto funcionamiento de la solicitud de registro,
y luego su posterior activac\'on. Para ello se cuenta con un usuario que tiene permisos de administraci\'on,
es decir, como se mencion\'on al hablar de los diagramas UML, el actor Admin, es quien puede
realizar las operaciones CRUD sobre los usarios. La Figure~\ref{formularioRegistro}, muestra
como se ve un formulario de registro en DermaUH 3.0 \newpage
\begin{figure}[ht] 
    \centering
    \includegraphics[width=0.6\textwidth]{Graphics/registro.png}
    \caption{Formulario de Registro} 
    \label{formularioRegistro}
\end{figure}

Una vez relleno el formulario, la solicitud es enviada al administrador, quien recibe un correo
alertando que tiene un nuevo usuario pendiente de respuesta. Acto seguido el administrador entra
al sistema con las credenciales correspondientes. La Figura~\ref{adminPage} muestra como se ve
la p\'agina principal de los usarios, tiene un contador de solicitudes pendientes, doctores activos
y residentes activos. Al presionar el bot\'on ``Pendientes'', se muestra la lista de usuarios que 
esperan confirmaci\'on para entrar al sistema, Figura~\ref{pendientes}. 
\begin{figure}[ht] 
    \centering
    \includegraphics[width=0.8\textwidth]{Graphics/adminPage.png}
    \caption{P\'agina del Administrador} 
    \label{adminPage}
\end{figure}
\begin{figure}[ht] 
    \centering
    \includegraphics[width=0.8\textwidth]{Graphics/pendientes.png}
    \caption{Usuarios Pendientes} 
    \label{pendientes}
\end{figure}

Si se fija, en el listado de usuarios pendientes, cada tarjeta posee un bot\'on ``Activar'',
al presionarlo se muestra el formulario de la Figure~\ref{formularioActivar}, el cual decide
que permisos tendr\'a el usuario en cuesti\'on. Solo se puede tener un permiso a la vez,
el front se configur\'o de tal forma que solo se puede escoger uno.\newpage

\begin{figure}[ht] 
    \centering
    \includegraphics[width=0.6\textwidth]{Graphics/activarFormulario.png}
    \caption{Formulario Activar Usuario} 
    \label{formularioActivar}
\end{figure}

Una vez escogido el rol del usuario, para este caso, se seleccion\'o doctor, se pulsa el bot\'on
activar. Si no hay ning\'un problema, deber\'ia aparecer un mensaje de usuario activado correctamente
como se muestra en la Figure~\ref{usuarioActivo}. Una vez activado el usuario, este recibe un correo 
electr\'onico v\'ia gmail de que puede utilizar los servicios de DermaUH 3.0.

\begin{figure}[ht] 
    \centering
    \includegraphics[width=0.6\textwidth]{Graphics/seActivoBien.png}
    \caption{Mensaje de que el usuario fue activado con \'exito} 
    \label{usuarioActivo}
\end{figure}

\section{Pruebas de Flujo Doctor}
Una vez analizadas las pruebas para el actor Admin, se prodecer\'a a realizar pruebas funcionales
con un usuario cuyo rol es doctor. Las pruebas que se llevar\'an a cabo ser\'an: CRUD sobre la entidad
paciente, creaci\'on de una lesi\'on, creaci\'on de un diagn\'ostico y subida de imagen, finalmente 
se compartir\'a la lesi\'on con otro usuario, se entrar\'a al sistema con sus credenciales y se
verificar\'a que la lesi\'on fue compartida exitosamente.

\subsection{CRUD sobre Paciente}
Esta prueba se realiza con el objetivo de comprobar si la aplicaci\'on soluciona de forma correcta
el RF1. Para ello se inciar\'a sesi\'on con un usuario con permisos de Doctor, como se muestra en la 
